\documentclass[10pt,a4paper]{amsart}
\usepackage[margin=19mm]{geometry}
\usepackage{amsmath,amssymb,mathtools}
\usepackage[scr=rsfs,cal=euler]{mathalfa}
\usepackage{xfrac}
\usepackage[unicode,hidelinks,bookmarks=false]{hyperref}
\newcommand{\ie}{i.e.\ }
\newcommand{\cf}{cf.\ }
\newcommand{\resp}{respectively\ }
\newcommand{\Subsection}{\subsection}
\providecommand{\nobreakdash}{\nobreak\hbox{-}}
\setlength{\emergencystretch}{2em}
\multlinegap0pt
\usepackage{amssymb}
\usepackage{mathtools}
\usepackage[shortlabels]{enumitem}
\usepackage{xcolor}
\usepackage{csquotes}
\usepackage{tikz}
\newtheorem{lemma}{Lemma}
\usepackage{cleveref}
\crefname{lemma}{Lemma}{Lemmas}
\newcommand{\N}{\mathbb{N}}
\title{Cubulation of Bruhat graphs}
\author{Alex Bishop, Elizabeth Mili\'cevi\'c, Anne Thomas}
\date{Source: arXiv:2504.03046v2 (CC BY 4.0)}
\begin{document}
\maketitle

\noindent\emph{Source and licence.} Cubulation of Bruhat graphs. Version arXiv:2504.03046v2 (CC BY 4.0), distributed by arXiv under the Creative Commons Attribution 4.0 International licence, http://creativecommons.org/licenses/by/4.0/, as recorded in the arXiv licence metadata for this exact version and shown on the abstract page https://arxiv.org/abs/2504.03046v2. Full licence notice: \texttt{licenses/arxiv-2504.03046-CC-BY-4.0.txt}.\par\smallskip

\noindent\textit{Editorial scope.} Counted unit: the article's lemma lem:limit-of-quantum (source0.tex lines 1101--1122). The article's complete original proof of it is delivered with it (source0.tex lines 1124--1204, one \texttt{proof} environment of 81 lines). No mathematics is written, repaired or re-derived: the statement and the proof are the article's own lines, in the article's own notation, and no retained line is substituted or re-flowed.  Retained uncounted as the source-local context the proof needs: lines 536--539.  Nothing was omitted from any retained span, so the package carries no omission record.  No retained line contains a citation, so the wrapper carries no bibliography and no bibliography entry.  No retained line contains a figure environment, an image-inclusion command or drawing code, so no image is delivered and the package declares no asset dependency.  Editorial formatting only, in addition: an AMS \texttt{amsart} preamble that restates the article's own declarations and macros used by the retained lines, namely the environments \texttt{lemma} on separate counters, together with the standard packages those lines need.  The retained environments print in this document as Lemma 1 in order of appearance, whereas the article prints the same items with the numbering of its own full document.  The complete native source of the article as distributed in the arXiv e-print for this exact version is bundled unchanged as \texttt{sources/175945/source0.tex}.
\begin{lemma}\label{lem:power-series-product-agreement}
	Let $F(q)$, $G(q)$, and $H(q)$ be power series, and define \[ \Phi(q) = F(q)H(q) \quad \mbox{and} \quad \Gamma(q) = G(q)H(q).\]
	Then for all $k \in \N$, if $F[k](q) = G[k](q)$, we have $\Phi[k](q) = \Gamma[k](q)$. 
\end{lemma}

\begin{lemma}\label{lem:limit-of-quantum}
  Let $n \geq 1$, and let $(g_j(q))_{j=0}^\infty$ be a sequence of polynomials of the form
	\[
	g_j(q)
	=
	(1 - q^{a_{1,j}})
	(1 - q^{a_{2,j}})
	\cdots
	(1 - q^{a_{n,j}}),
	\]
	where for all $1 \leq i \leq n$ and all $j \geq 0$, the $a_{i,j}$ are integers satisfying \[ 1 \leq a_{1,j} \leq a_{2,j} \leq \cdots \leq a_{n,j}.\]
	Suppose that $F(q)$ is a power series such that for each $j \geq 0$, the Taylor polynomials $F[j](q)$ and $g_j[j](q)$ satisfy \[ F[j](q) = g_j[j](q).\] 
	Then either $F(q) = 1$, or there exists a positive integer $M$ and an integer $N$ with $1 \leq N \leq n$, such that 
	\[
	F(q) =
	(1 - q^{a_{1,M}})
	(1 - q^{a_{2,M}})
	\cdots
	(1 - q^{a_{N,M}})
	\]
	and for all $j \geq M$, we have $g_j(q) = g_M(q)$.
\end{lemma}

\begin{proof}
	We will prove the statement by induction on $n$. 
	Suppose $n = 1$. Then for each $j \geq 1$, we have $
	g_j(q) = 1 - q^{a_{1,j}}$.
	If $F(q) \neq 1$, then there must be some $m \in \N$ so that $g_m[m](q) \neq 1$. Then \[ g_m[m](q) = 1 - q^{a_{1,m}},\] and so in particular, $a_{1,m} \leq m$.
	Now for all $j \geq m$, we have $j \geq a_{1,m}$, so
	\[
	g_j[j](q) = 1 - q^{a_{1,m}}
	\]
	as well. Hence $g_j(q) = g_m(q)$ for every $j \geq m$. But then for every $j \geq m$, we also have
	\[
	F[j](q) = 1 - q^{a_{1,m}},
	\]
	and hence $F(q) = 1 - q^{a_{1,m}}$. Put $M = m$ and $N = n = 1$, and we have established the result for $n = 1$.
	
	For the inductive step, we again have that if $F(q) \neq 1$ then $g_m[m](q) \neq 1$ for some $m \in \N$. Now $a_{1,m} \leq \dots \leq a_{n,m}$, so if we expand out the product $g_m(q)$ we obtain
	\[
	g_m(q) = 1 - c_m q^{a_{1,m}} + \mbox{higher degree terms},
	\]
	where the coefficient $c_m \neq 0$ is the number of exponents $a_{1,m},\dots,a_{n,m}$ to be equal to $a_{1,m}$. We note that $a_{1,m} \leq m$. Hence for all $j \geq m$, since $j \geq a_{1,m}$ we have
	\[
	g_j[j](q) = 1 - c_m q^{a_{1,m}} + \mbox{higher degree terms}.
	\]
	Thus in particular, since $a_{1,j} \leq \dots \leq a_{n,j}$, we have $a_{1,j} = a_{1,m}$ for all $j \geq m$.
	So for all $j \geq m$ we have 
	\[
	\frac{g_j(q)}{1-q^{a_{1,m}}} =
	(1 - q^{a_{2,j}})
	\cdots
	(1 - q^{a_{n,j}}).
	\]
	We can thus define a sequence of polynomials $(h_j(q))_{j=0}^\infty$ by
	\[
	h_j(q) = \frac{g_{m+j}(q)}{1-q^{a_{1,m}}} =
	(1 - q^{a_{2,m+j}})
	\cdots
	(1 - q^{a_{n,m+j}}).
	\]
	
	Now as the quotient $1/(1-q^{a_{1,m}})$ 
	is itself a power series, we can view each $h_j(q)$ as the product of the polynomial $g_{m+j}(q)$ with this power series, and we can also define the product 
	\[
	G(q) = F(q) \left(
	\frac{1}{1-q^{a_{1,m}}} \right) = \frac{F(q)}{1-q^{a_{1,m}}}.
	\]
	Since $F[j](q) = g_j[j](q)$ for all $j \geq 0$, we can thus apply \cref{lem:power-series-product-agreement} to see that for all $j \geq 0$, 
	\begin{eqnarray*}
	G[m+j](q)  & = & \left( 
	\frac{F(q)}{1-q^{a_{1,m}}} 
	\right)
	[m+j](q) \\
	& = &
	\left( 
	\frac{g_{m+j}(q)}{1-q^{a_{1,m}}}
	\right)
	[m+j](q) \\ & = & h_j[m+j](q).
	\end{eqnarray*}
	Therefore $G[j](q) = h_j[j](q)$ for all $j \geq 0$. 
	
	
	By inductive assumption, since each $h_j(q)$ is a product of $(n-1)$ factors, either $G(q) = 1$, or there is an $M' \geq 1$ and an integer $N'$ with $1 \leq N' -1 \leq n-1$, such that $G(q)$ is the product of $(N' - 1)$ factors as follows:
	\[
	G(q) 
	=
	(1 - q^{a_{2,M'}})
	\cdots
	(1 - q^{a_{N',M'}}),
	\]
	and for all $j \geq M'$, we have $h_j(q) = h_{M'}(q)$. Recall also from above that $a_{1,j} = a_{1,m}$ for each $j \geq m$. 
	
	If $G(q) = 1$, then $F(q) = 1 - q^{a_{1,m}}$ and $g_j(q) = g_m(q)$ for all $j \geq m$, and we are done with $M = m$ and $N = 1 \leq n$. Now assume $G(q) \neq 1$. Then we have in particular that $a_{i,m+j} = a_{i,m+M'}$ for all $2 \leq i \leq N'$ and all $j \geq M'$.  Put $M = m + M'$. Then $a_{1,m} = a_{1,M}$ and $a_{i,M'} = a_{i,M}$ for all $2 \leq i \leq n$, and so $F(q) = (1 - q^{a_{1,m}})G(q)$ is the product of $1 + (N'-1) = N'$ factors as follows:
	\[
	F(q)
	=
	(1-q^{a_{1,M}})
	(1 - q^{a_{2,M}})
	\cdots
	(1 - q^{a_{N',M}}),
	\]
	and for each $j \geq M$, we have $g_j(q) = g_M(q)$. Letting $N = N'$, this completes the proof of the inductive step.
\end{proof}

\end{document}
